\section{Methodology}

\subsection{Framework Overview}

We propose a six-layer hierarchical framework that decomposes literary quality into increasingly interpretive dimensions. The framework distinguishes between lower layers (L1-L3) that assess objective textual properties through rule-based computational metrics and upper layers (L4-L6) that evaluate interpretive literary qualities through LLM-based critical analysis. This decomposition reflects an empirically testable hypothesis: formal textual features and interpretive literary qualities constitute largely independent dimensions requiring different analytical approaches.

Figure~\ref{fig:framework} illustrates the overall architecture, and Table~\ref{tab:framework_layers} summarizes the six-layer hierarchy. The modular design permits selective evaluation using layer subsets depending on analytical objectives. Our experimental validation focuses on Layers 4, 5, and 6, which span cultural, emotional-psychological, and existential-philosophical dimensions with measurable criteria grounded in established theoretical frameworks.

\begin{figure*}[t]
\centering
\includegraphics[width=\textwidth]{sections/figures/framework_diagram.png}
\caption{SAGE framework architecture. (a)~Six-layer hierarchy: Layers 1--3 assess textual properties through rule-based metrics (dashed); Layers 4--6 evaluate interpretive dimensions through LLM-based analysis (solid). (b)~Evaluation pipeline: each story is assessed in two modes (content-limit and title-limit) across three layers via a dual-track architecture comprising a five-round iterative evaluator and an independent validator, with inter-rater agreement analysis.}
\label{fig:framework}
\end{figure*}

\begin{table}[t]
\centering
\small
\resizebox{\columnwidth}{!}{%
\begin{tabular}{clp{7cm}}
\hline
\textbf{Layer} & \textbf{Focus} & \textbf{Key Dimensions} \\
\hline
\multicolumn{3}{l}{\textit{Rule-Based Evaluation (Computational Metrics)}} \\
\hline
L1 & Lexical Quality & Vocabulary richness, syntactic complexity, rhetorical density \\
L2 & Narrative Structure & Entity coherence, emotional arc trajectories, event sequence patterns \\
L3 & Thematic Content & Topic diversity, semantic networks, conceptual structure \\
\hline
\multicolumn{3}{l}{\textit{LLM-Based Evaluation (Interpretive Analysis)}} \\
\hline
L4 & Cultural & \textbf{IPD}: Intersectional Power Dynamics; \\
 & Representation & \textbf{CVP}: Cultural Voice \& Perspective; \\
 & & \textbf{CSP}: Cultural Specificity; \\
 & & \textbf{CPC}: Cultural Pattern Complexity \\
\hline
L5 & Emotional- & \textbf{AC}: Affective Complexity; \\
 & Psychological & \textbf{PI}: Psychological Interiority; \\
 & & \textbf{EG}: Emotional Granularity; \\
 & & \textbf{ENC}: Emotional-Narrative Coherence \\
\hline
L6 & Existential- & \textbf{LP}: Life Philosophy; \\
 & Philosophical & \textbf{MR}: Moral Reflection; \\
 & & \textbf{HC}: Human Condition; \\
 & & \textbf{ME}: Meaning Exploration \\
\hline
\end{tabular}
}
\caption{Six-layer hierarchical evaluation framework}
\label{tab:framework_layers}
\end{table}

\subsection{Rule-Based Evaluation Layers (L1-L3)}

The lower three layers employ computational metrics to assess textual properties amenable to algorithmic measurement, including vocabulary richness, syntactic complexity, entity coherence, emotional arc trajectories, topic diversity, and semantic network structure. These layers provide efficient, reproducible assessment of surface and discourse-level features. As the present study focuses on validating the interpretive layers (L4--L6), the rule-based layers are summarized in Table~\ref{tab:framework_layers} and their experimental validation is left for future work.

\subsection{LLM-Based Evaluation Layers (L4-L6)}

The upper three layers employ large language models to assess interpretive dimensions requiring cultural knowledge, emotional intelligence, and philosophical sophistication. All three layers form the focus of our experimental validation, each grounded in established scholarly traditions.

\subsubsection{Layer 4: Cultural Representation}

Layer 4 evaluates how texts represent culture, power structures, and social relations. The evaluation draws on multiple theoretical traditions including Bourdieu's field theory and cultural capital, Said's critique of Orientalism, Spivak's work on subaltern representation, Geertz's concept of thick description, and structural anthropology. These frameworks provide explicit criteria to assess cultural representation quality independently of emotional or philosophical content, enabling systematic evaluation of how texts engage with social reality as structured phenomenon.

Building on this theoretical foundation, we define four dimensions that collectively characterize cultural representation quality:

\textbf{Intersectional Power Dynamics (IPD)} examines how power operates through social structures and relationships. Grounded in Bourdieu's theory of cultural capital alongside Marxist, feminist, and postcolonial frameworks, this dimension assesses the complexity and configuration of power relations enacted through various social axes including class, gender, race, kinship, occupation, age, religion, and institutional position. The evaluation prioritizes structural complexity and textual evidence over mere presence of identity markers, recognizing that power may manifest through single or multiple axes depending on narrative context. High scores require concrete textual anchors demonstrating sophisticated treatment of power asymmetries, symbolic violence, or structural constraints rather than superficial acknowledgment of social categories.

\textbf{Cultural Voice and Perspective (CVP)} analyzes narrative positioning relative to cultural material. Drawing on Said's critique of Orientalism and Spivak's work on subaltern representation, this dimension distinguishes insider perspectives demonstrating cultural authority from outsider gazes marked by exoticization or ethnographic distance. Textual markers including linguistic register, deployment of untranslated terms, focalization strategies, and treatment of cultural knowledge reveal whether narrative voice presents cultural systems coherently from within or frames them as spectacle from without. The dimension operates independently of authorial identity, focusing instead on how narrative voice positions itself epistemically toward represented cultures.

\textbf{Cultural Specificity (CSP)} quantifies density and precision of cultural particularities across temporal, spatial, social, and material domains. Informed by Geertz's concept of thick description, this dimension assesses how texts ground representations in concrete details of historical period, geographic location, social customs, and material culture. High specificity manifests through precise temporal markers (not merely "Victorian era" but specific years or events), detailed geographic localization, attention to ritual and practice, and deployment of material objects embedding cultural meaning. This dimension captures representational thickness independently of whether details serve emotional engagement or philosophical argumentation.

\textbf{Cultural Pattern Complexity (CPC)} evaluates structural sophistication of cultural patterns through synthesis of structural anthropology and cultural systems theory. The dimension assesses: (1) whether social relationships exhibit multi-layered structural logic or reduce to binary oppositions; (2) whether narrative conflicts reflect institutional and systemic forces or rely on personalized antagonisms; (3) whether texts demonstrate critical reflexivity toward represented cultural systems or present them with narrative transparency; (4) whether cultural patterns exhibit archetypal depth or rely on genre clichés and formulaic conventions. Together these criteria characterize whether cultural representation achieves the structural sophistication distinguishing literary achievement from formulaic production.

\subsubsection{Layer 5: Emotional-Psychological Representation}

Layer 5 evaluates how texts represent emotions, psychological states, and affective experiences. The evaluation draws on affect theory (Sedgwick, Berlant, Ngai), narrative psychology (Cohn, Bruner, Bakhtin), emotion granularity research (Barrett), and New Critical concepts of organic unity and motivated emotion (James Wood). These frameworks provide systematic criteria to assess emotional-psychological representation quality independently of cultural power structures or existential-philosophical themes.

Building on this theoretical foundation, we define four dimensions that collectively characterize emotional-psychological representation quality:

\textbf{Affective Complexity (AC)} measures multiplicity, contradiction, and evolution of emotional states. Grounded in affect theory (Sedgwick, Berlant, Ngai), this dimension evaluates whether emotions appear as simple singular affects or exhibit multi-layered, contradictory, and evolving patterns. High complexity manifests through coexisting contradictory emotions, affective ambiguity, transformation across narrative progression, and avoidance of emotional binaries. The dimension focuses on structural properties of emotional representation rather than intensity or explicitness of expression.

\textbf{Psychological Interiority (PI)} assesses depth and accessibility of characters' inner psychological worlds. Drawing on narrative psychology (Cohn, Bruner, Bakhtin), this dimension evaluates the degree to which narrative provides access to thoughts, motivations, self-awareness, and mental processes beyond observable emotions. High interiority manifests through interior monologue, free indirect discourse, psychological focalization, or other techniques revealing characters' reasoning, memory, fantasy, and conflicting desires. The dimension distinguishes psychological depth from mere behavioral description.

\textbf{Emotional Granularity (EG)} evaluates precision and differentiation of emotional vocabulary and concepts. Grounded in emotion granularity research (Barrett) and affective linguistics, this dimension assesses whether emotions receive fine-grained specification or remain vague and undifferentiated. High granularity manifests through precise emotion terms distinguishing similar states (melancholic vs. wistful vs. despondent) or through detailed behavioral and physiological descriptions that specify emotional experiences without generic labels. The dimension captures linguistic precision in emotional representation.

\textbf{Emotional-Narrative Coherence (ENC)} evaluates whether emotional moments arise organically from narrative situation and character or appear unmotivated and manipulative. Grounded in New Critical concepts of organic unity and James Wood's notion of motivated emotion, this dimension assesses logical-causal relationships between emotions and their narrative contexts. High coherence requires that emotional responses possess clear antecedent motivation in plot or character development, exhibit proportionality to narrative stakes, and serve narrative functions rather than appearing imposed for effect. The dimension identifies sentimentality, melodrama, and unearned pathos as markers of low coherence.

\subsubsection{Layer 6: Existential-Philosophical Representation}

Layer 6 evaluates philosophical depth and engagement with existential themes addressing fundamental questions of human existence. The evaluation draws on existentialist philosophy (Heidegger, Sartre, Camus), moral philosophy (Levinas, MacIntyre), philosophical anthropology (Arendt, Nussbaum), and hermeneutic traditions (Ricoeur, Gadamer). These frameworks provide systematic criteria to assess existential-philosophical representation quality independently of cultural power analysis or emotional-psychological complexity.

Building on this theoretical foundation, we define four dimensions that collectively characterize existential-philosophical representation quality:

\textbf{Life Philosophy (LP)} examines engagement with fundamental questions of human existence, authenticity, freedom, and life's meaning. Drawing on existentialist philosophy (Heidegger's Being-towards-death, Sartre's radical freedom, Camus's Absurdism) alongside Eastern philosophical traditions including Confucianism, Daoism, and Buddhism, this dimension assesses whether texts grapple substantively with questions of how to live, what constitutes authentic existence, and what gives life meaning. High scores require sustained philosophical inquiry through narrative situation and character confrontation with existential questions, not merely surface-level philosophical gestures or generic wisdom statements. The dimension distinguishes genuine philosophical exploration, where characters or narrative structure embody philosophical positions through lived experience, from didactic insertion of philosophical quotations or abstract discussions disconnected from narrative embodiment.

\textbf{Moral Reflection (MR)} assesses depth of ethical inquiry, moral complexity, and engagement with questions of right conduct and moral responsibility. Grounded in moral philosophy (Levinas's ethics of alterity, MacIntyre's virtue ethics, Confucian ethics of reciprocity and ritual propriety), this dimension evaluates whether texts present moral dilemmas with genuine complexity and ambiguity or reduce ethical questions to clear-cut binaries. High moral reflection manifests through situations where competing moral frameworks yield conflicting obligations, where characters confront irreducible moral ambiguity without simple resolution, or where narrative structure itself questions conventional moral certainties. The dimension captures sophistication of ethical thinking rather than moral content, recognizing that morally troubling or transgressive texts may demonstrate profound moral reflection through their questioning of ethical assumptions.

\textbf{Human Condition (HC)} evaluates exploration of universal aspects of human experience including mortality, suffering, finitude, vulnerability, and the search for meaning in face of existential limits. Based on philosophical anthropology (Arendt's analysis of human plurality and natality, Nussbaum's capabilities approach highlighting human vulnerability and interdependence), this dimension assesses whether texts engage meaningfully with constraints and possibilities defining human existence. High scores require concrete narrative engagement with existential realities such as aging, death, loss, bodily limitation, and temporal finitude, treated not as plot devices but as fundamental conditions shaping human experience and meaning-making. The dimension distinguishes texts that confront the human condition with philosophical seriousness from those employing mortality or suffering merely for emotional effect or narrative convenience.

\textbf{Meaning Exploration (ME)} examines how narratives construct, question, or subvert frameworks of meaning and interpretive horizons. Drawing on hermeneutic philosophy (Ricoeur's narrative identity and productive interpretation, Gadamer's fusion of horizons), this dimension assesses whether texts thematize the construction of meaning itself: how individuals and communities create coherent narratives from experience, how interpretive frameworks shape understanding, and how meaning proves fragile or contested. High meaning exploration manifests through narrative self-reflexivity about meaning-making processes, characters whose primary struggle involves interpretation and sense-making, or structural techniques that foreground the contingency of narrative coherence. The dimension captures meta-level engagement with hermeneutic processes rather than mere presence of symbolic meaning or thematic significance.

\subsection{Dual-Mode Evaluation Framework}

To probe different aspects of LLM evaluation capacity, we implement two evaluation modes that vary the information provided to the model.

\textbf{Content-Limit Mode} provides only the story text without title, author, or metadata. This mode requires the LLM to assess literary qualities based solely on observable textual evidence without relying on memorized information about specific works or authors. Content-limit evaluation tests whether LLMs can apply generalizable critical frameworks to unfamiliar texts, analyzing cultural representation, emotional complexity, and other interpretive dimensions from textual patterns alone. Prompts in this mode explicitly instruct the model to avoid using prior knowledge and cite specific passages to ground all assessments.

\textbf{Title-Limit Mode} provides only the work's title and author without story text. This mode requires the LLM to synthesize scholarly consensus and critical reception from its training corpus. Title-limit evaluation tests whether LLMs can accurately recall and integrate expert literary criticism to produce meta-critical assessments. Prompts in this mode instruct the model to draw on academic analyses, cite specific critics or scholarly perspectives where relevant, and distinguish reliable critical consensus from speculative interpretations.

Comparing content-limit and title-limit modes reveals whether LLM evaluations derive from textual analysis, memorized critical opinions, or a combination of both sources. Systematic differences between modes indicate reliance on prior knowledge, while high agreement suggests that LLMs can perform both textual analysis and critical synthesis yielding consistent assessments.

\subsection{Dual-Track Evaluation Architecture}

To ensure reliable and interpretable assessments, we employ two evaluators with complementary roles: an iterative evaluator that performs multi-round self-reflection and an independent validator that provides cross-verification.

\subsubsection{Iterative Evaluator}

The iterative evaluator performs assessment across five structured rounds, progressively refining judgments while tracking confidence and reasoning quality. Each round addresses specific sources of bias or error through targeted self-reflection prompts.

\textbf{Round 1 (Initial Assessment):} The model receives the story (content-limit) or title/author (title-limit) along with theoretical definitions for each dimension. The model extracts relevant content, provides initial scores for all dimensions, assigns confidence ratings, and justifies assessments with specific textual evidence or critical citations.

\textbf{Round 2 (Projection Bias Check):} The model reviews Round 1 output and examines whether it imposed inappropriate assumptions about emotional expression norms, cultural contexts, or aesthetic conventions. The model revises scores if projection bias is detected or confirms Round 1 assessments if no bias is found.

\textbf{Round 3 (Layer Boundary Compliance):} The model verifies that it evaluated only the target layer's dimensions without conflating distinct layers. For Layer 4, this involves checking that cultural assessments did not incorporate emotional or philosophical considerations. For Layer 5, this involves verifying that emotional assessments did not evaluate cultural power structures or existential themes.

\textbf{Round 4 (Evidence Sufficiency):} The model examines the textual evidence or critical citations grounding each dimension score. The model adjusts confidence ratings (not scores) if evidence proves ambiguous or sparse, explicitly noting where inferences extend beyond clear textual support.

\textbf{Round 5 (Final Consolidation):} The model provides final scores with consolidated reasoning, summarizing key evidence, noting which dimensions exhibited convergence (minimal score change from Round 4 to Round 5 indicates stable assessment), and assigning overall evaluation confidence.

All rounds produce structured JSON outputs specifying dimension scores, confidence ratings, reasoning, and round-specific reflection data. This multi-round architecture addresses known LLM evaluation inconsistency by providing explicit opportunities to correct biases and verify evidence grounding.

\subsubsection{Independent Validator}

The independent validator evaluates the same texts without access to iterative evaluator outputs, providing cross-verification and hallucination detection. The validator performs single-pass assessment producing:

\begin{enumerate}
\item \textbf{Independent Scores:} Dimension-level scores and confidence ratings based solely on the validator's analysis.

\item \textbf{Validation Review:} Critical examination of whether the iterative evaluator's reasoning exhibits projection bias, hallucination, layer boundary violations, or unsupported inferences. The review identifies specific issues with severity ratings.

\item \textbf{Agreement Analysis:} Comparison between validator scores and iterative evaluator final scores, noting dimensions with strong agreement vs. significant discrepancies.

\item \textbf{Trust Level:} Overall qualitative assessment indicating confidence in the iterative evaluation reliability.
\end{enumerate}

Measuring agreement between iterative evaluator and independent validator provides inter-rater reliability evidence. High agreement suggests that both evaluators converge on similar interpretations when applying the same theoretical frameworks. Systematic disagreements on specific dimensions or text types reveal where LLM evaluation proves less reliable.

